\documentclass[11pt]{article}
\usepackage[margin=1in]{geometry}
\usepackage{amsmath,amssymb,amsthm,hyperref}
\newtheorem{theorem}{Theorem}
\title{An additive-game counterexample to the core vertex-count law}
\author{ziangni-sys}
\date{Conjecture 00000008232}
\begin{document}
\sloppy
\maketitle
\noindent Both source versions assert that the number of vertices of the core
of a convex geometry (``anti-convex games'') is $2^{n-1}$. We disprove this
vertex-count clause using a two-player additive game on the full Boolean
convex geometry. We verify both conventional payoff and cost core inequalities.

\section*{The actual geometry and game}
Let $N=\{0,1\}$ and take all subsets of $N$ as feasible coalitions.
The closure operator is $\operatorname{cl}(S)=S$. This is a finite convex
geometry: closure is extensive, monotone and idempotent, and
$\operatorname{cl}(\varnothing)=\varnothing$. Its anti-exchange axiom says
that for distinct $x,y$ outside $\operatorname{cl}(S)$, if
$x\in\operatorname{cl}(S\cup\{y\})$, then
$y\notin\operatorname{cl}(S\cup\{x\})$. Here the antecedent is impossible:
$x\in S\cup\{y\}$ would imply $x\in S$ or $x=y$.
Every coalition is closed and feasible. Removing any element from a
nonempty feasible coalition again gives a feasible coalition. Thus the full
Boolean family also satisfies the accessible feasible-family convention.

Define the normalized game $v(S)=|S|$. The finite-set cardinality identity
\[
 |S\cup T|+|S\cap T|=|S|+|T|
\]
shows that $v$ is modular. In particular it is both supermodular (the usual
convex-game inequality) and submodular (the reversed inequality used for
concave or cost games). All these inequalities hold with equality.

\begin{theorem}
For this game, under either payoff or cost conventions, the core is
$\{(1,1)\}$ and has one vertex. The asserted vertex count is
$2^{2-1}=2$, so the conjecture is false.
\end{theorem}
\begin{proof}
The payoff core consists of $x=(x_0,x_1)\in\mathbb R^2$ such that
\[
 x_0+x_1=v(N)=2,\qquad \sum_{i\in S}x_i\ge v(S)
 \quad\text{for every feasible }S\subseteq N.
\]
The singleton coalitions require $x_0\ge1$ and $x_1\ge1$.
Efficiency forces $x_0=x_1=1$. Conversely, $(1,1)$ satisfies every coalition
inequality with equality. Hence the entire payoff core is $\{(1,1)\}$.
For the cost core the inequalities are reversed: $x_0\le1$, $x_1\le1$,
with the same efficiency equality. The same argument gives exactly the
same singleton core.

A point of a convex set is extreme if it cannot be written as a strict
convex combination of two distinct points of the set. The unique point of
a singleton is extreme, and there are no other points. Both cores are
therefore genuine singleton polytopes with exactly one vertex, contrary to
$2^{n-1}=2$ for $n=2$.
\end{proof}
\newpage
\section*{Source conventions and scope}
The source's term ``anti-convex games'' is not separately defined in either
language. Our example satisfies both standard non-strict game inequalities,
and its feasible family is an actual convex geometry. No strict
nonmodularity condition appears in the source. We do not add such a
condition or claim a result about a different restricted game class.
The example also avoids ambiguity about payoff versus cost conventions by
computing both entire cores.

The full Boolean family is a permitted convex geometry, including the
identity closure as its simplest case. The restriction to feasible
coalitions introduces no omissions in this example, since all coalitions
are feasible. Additive games are normalized ($v(\varnothing)=0$) and meet
the usual convex and concave game conditions. A core need not have full
dimension to be a polytope with vertices: a singleton is a zero-dimensional
polytope with one vertex.

The vertex-count assertion is a conjunct of the original claim. Refuting
it suffices to disprove the conjecture; no claim about the other extreme-ray,
chain-count or entropy assertions is required.

\section*{Actual Lean objects and checked theorems}
The player type is \texttt{Fin 2}, coalitions are actual finite subsets,
and allocations are real-valued functions on the players.
\texttt{IsConvexGeometry} states the closure axioms and anti-exchange
condition, proved by \texttt{boolean\_convex\_geometry} for the identity
closure. \texttt{every\_coalition\_feasible} and
\texttt{feasible\_accessible} verify the feasible-family correspondence.

The actual game value \texttt{game s} is the real cast of
\texttt{s.card}. \texttt{game\_modular} proves the cardinality identity;
\texttt{game\_supermodular} and \texttt{game\_submodular} prove both
standard inequalities. \texttt{payoffCore} and \texttt{costCore} each
impose the efficiency equality and every feasible coalition inequality.
They do not assume that only singleton constraints matter.
\texttt{payoff\_core\_singleton} and \texttt{cost\_core\_singleton}
prove equality of the full cores with the singleton allocation.

The formalization uses Mathlib's actual \texttt{Set.extremePoints} over
$\mathbb R$, whose definition uses membership in strict open line segments.
Its proved singleton theorem computes both extreme-point sets.
\texttt{vertex\_count\_fails} computes their actual \texttt{Set.ncard}
as $1$ and proves that neither is $2^{\#N-1}$.

\section*{Reproduction}
Lean 4.19.0 and Mathlib are pinned at commit
\begin{center}\small\texttt{c44e0c8ee63ca166450922a373c7409c5d26b00b}.\end{center}
From \texttt{lean/}, run \texttt{lake update}, \texttt{lake exe cache get},
and \texttt{lake build}; then check the whole source with
\begin{center}\small\texttt{lake env lean -DwarningAsError=true Main.lean}.\end{center}
The axiom audits use only standard logical axioms. No statements are
admitted, and no custom axioms or native decision procedures are used.
\end{document}
